\chapter{Pseudodifferential operators on Euclidean space}\label{ch:euclidean}
\providecommand{\Amp}[2]{\mathcal{A}^{#1}_{#2}}% amplitudes of order #1 and weight #2

This chapter constructs the basic algebra of pseudodifferential operators on
$\R^n$, the ``traditional'' calculus with symbols of type $(1,0)$ whose
coefficients are uniformly smooth in the base variables. The starting point is
the Fourier representation of a differential operator with coefficients in
$\CIb(\R^n)$,
\begin{equation*}
  P(x,D)u(x)=\int e^{i(x-y)\cdot\xi}P(x,\xi)u(y)\,dy\,\dbar\xi
  =\int e^{ix\cdot\xi}P(x,\xi)\hat u(\xi)\,\dbar\xi,\qquad u\in\cS(\R^n),
\end{equation*}
where $P(x,\xi)=\sum_{\abs\alpha\le m}p_\alpha(x)\xi^\alpha$ is the
characteristic polynomial. We generalize this formula by replacing the
polynomial $P(x,\xi)$ by a \emph{symbol}, a smooth function satisfying
uniform estimates at infinity, and we even allow it to depend on the
``incoming'' variable $y$, so that operators are given by
\begin{equation}\label{euclidean:eq:amplitude-operator}
  Au(x)=\int e^{i(x-y)\cdot\xi}a(x,y,\xi)u(y)\,dy\,\dbar\xi,\qquad u\in\cS(\R^n).
\end{equation}
Such integrals are in general not absolutely convergent.
\Cref{euclidean:sec:symbols} develops the symbol spaces, including
asymptotic summation. In \cref{euclidean:sec:operators} the integral
\eqref{euclidean:eq:amplitude-operator} is interpreted as an oscillatory
integral and shown to define a continuous operator on $\cS(\R^n)$.
\Cref{euclidean:sec:reduction} shows that every such operator can be written
with an amplitude independent of $y$ (reduction to the left symbol), and
\cref{euclidean:sec:composition} deduces from this that the operators form
an order-filtered $*$-algebra. The remaining sections develop the
consequences: principal symbols, ellipticity and parametrices, elliptic
regularity, boundedness on $L^2$ and on (weighted) Sobolev spaces,
polyhomogeneous (classical) operators and invariance under affine changes of
variables. Coordinate invariance, and hence operators on manifolds, is the
subject of \cref{ch:microlocal,ch:manifolds}.

Throughout, $\ang{\xi}=(1+\abs\xi^2)^{1/2}$, $D=-i\pa$,
$\dbar\xi=(2\pi)^{-n}d\xi$, and the Schwartz space carries the norms
$\norm\phi_k=\max_{\abs\alpha+\abs\beta\le k}\sup_x\abs{x^\alpha D^\beta\phi(x)}$
of \cref{ch:distributions}. The bilinear pairing of a distribution $u$ with a
test function $\phi$ is written $u(\phi)$ or $\pair u\phi$, and
$\inner uv=\int u\bar v\,dx$ is the $L^2$ inner product.

% ======================================================================
\section{Symbols}\label{euclidean:sec:symbols}

A polynomial $p(\xi)$ of degree at most $m$ satisfies
$\abs{p(\xi)}\le C\ang\xi^m$, and since $D^\alpha_\xi p$ has degree at most
$m-\abs\alpha$, also $\abs{D^\alpha_\xi p(\xi)}\le C_\alpha\ang\xi^{m-\abs\alpha}$.
If $P(x,\xi)$ is the characteristic polynomial of a differential operator of
order $m$ with coefficients in $\CIb(\R^n)$ (all derivatives bounded) then
\begin{equation*}
  \abs{D^\alpha_xD^\beta_\xi P(x,\xi)}\le C_{\alpha\beta}\ang\xi^{m-\abs\beta}
  \quad\text{on }\R^n\times\R^n,\ \alpha,\beta\in\N_0^n.
\end{equation*}
There is no reason to have the same number of base and fibre variables, so in
general we define:

\begin{definition}\label{euclidean:def:symbols}
  For $m\in\R$ the space $S^m(\R^p;\R^n)$ of \emph{symbols of order $m$}
  consists of the functions $a\in\CI(\R^p\times\R^n)$ satisfying
  \begin{equation}\label{euclidean:eq:symbol-est}
    \abs{D^\alpha_zD^\beta_\xi a(z,\xi)}\le C_{\alpha\beta}\ang\xi^{m-\abs\beta}
    \quad\text{on }\R^p\times\R^n,\ \text{for all }\alpha\in\N_0^p,\ \beta\in\N_0^n.
  \end{equation}
  It is given the topology of the norms
  \begin{equation*}
    \norm a_{N,m}=\max_{\abs\alpha+\abs\beta\le N}\ \sup_{(z,\xi)\in\R^p\times\R^n}
    \ang\xi^{-m+\abs\beta}\abs{D^\alpha_zD^\beta_\xi a(z,\xi)},\qquad N\in\N_0.
  \end{equation*}
  We set $S^\infty(\R^p;\R^n)=\bigcup_mS^m(\R^p;\R^n)$ and
  $S^{-\infty}(\R^p;\R^n)=\bigcap_mS^m(\R^p;\R^n)$, the latter with the
  topology given by all the norms $\norm\cdot_{N,m}$, $N\in\N_0$, $m\in\R$;
  its elements are called \emph{residual} symbols.
\end{definition}

\begin{remark}\label{euclidean:rem:symbols-open}
  The estimates are uniform in the ``base'' variables $z$; this is the
  analogue of the condition that the coefficients be in $\CIb$. For an open
  set $\Omega\subset\R^p$ one defines $S^m(\Omega;\R^n)$ in the same way,
  requiring \eqref{euclidean:eq:symbol-est} on $\Omega\times\R^n$. The class
  of symbols is not the only reasonable one (see
  \cref{euclidean:ex:rho-delta}), but it is large enough to contain everything
  we need at first and small enough for the calculus to work easily.
\end{remark}

Since $\ang\xi^{m}\le\ang\xi^{m'}$ for $m\le m'$ we have
$\norm a_{N,m'}\le\norm a_{N,m}$, so there are continuous inclusions
$S^m(\R^p;\R^n)\hookrightarrow S^{m'}(\R^p;\R^n)$ for $m\le m'$; the symbol
spaces form a filtration of $S^\infty(\R^p;\R^n)$. A residual symbol is a
smooth function all of whose derivatives are bounded by
$C_{\alpha\beta N}\ang\xi^{-N}$ for every $N$.

We shall use repeatedly the following elementary inequality.

\begin{lemma}[Peetre's inequality]\label{euclidean:lem:peetre}
  For all $x,y\in\R^n$ and $s\in\R$,
  \begin{equation*}
    \ang{x+y}^s\le2^{\abs s/2}\ang x^s\ang y^{\abs s}.
  \end{equation*}
\end{lemma}
\begin{proof}
  Since $\abs{x+y}^2\le2\abs x^2+2\abs y^2$,
  $1+\abs{x+y}^2\le2(1+\abs x^2)(1+\abs y^2)$, i.e.\
  $\ang{x+y}\le\sqrt2\ang x\ang y$. This gives the case $s\ge0$. Applied to
  $x=(x+y)+(-y)$ it gives $\ang x\le\sqrt2\ang{x+y}\ang y$, i.e.\
  $\ang{x+y}^{-1}\le\sqrt2\ang x^{-1}\ang y$, and raising this to the power
  $\abs s$ gives the case $s<0$.
\end{proof}

\begin{lemma}\label{euclidean:lem:symbol-basic}
  Let $m,m'\in\R$.
  \begin{enumerate}[label=\textup{(\roman*)}]
  \item $\ang\xi^m\in S^m(\R^p;\R^n)$.
  \item Multiplication defines a continuous bilinear map
    $S^m(\R^p;\R^n)\times S^{m'}(\R^p;\R^n)\to S^{m+m'}(\R^p;\R^n)$.
  \item $D^\alpha_z$ maps $S^m(\R^p;\R^n)$ continuously into itself and
    $D^\beta_\xi$ maps it continuously into $S^{m-\abs\beta}(\R^p;\R^n)$.
  \item $S^m(\R^p;\R^n)$ is a Fr\'echet space.
  \item If $a\in\CI(\R^p\times\R^n)$ has all derivatives bounded and
    $a(z,\xi)=0$ for $\abs\xi\ge R$, then $a\in S^{-\infty}(\R^p;\R^n)$.
  \end{enumerate}
\end{lemma}
\begin{proof}
  (i) By induction on $\abs\beta$,
  $D^\beta_\xi\ang\xi^m=\ang\xi^{m-2\abs\beta}P_\beta(\xi)$ with $P_\beta$ a
  polynomial of degree at most $\abs\beta$: indeed
  $D_{\xi_j}(\ang\xi^{s}P)=\ang\xi^{s-2}\bigl(-is\xi_jP+\ang\xi^2D_{\xi_j}P\bigr)$
  and the polynomial in brackets has degree at most $\deg P+1$. Hence
  $\abs{D^\beta_\xi\ang\xi^m}\le C_\beta\ang\xi^{m-\abs\beta}$.

  (ii) By Leibniz' formula,
  \begin{equation*}
    \ang\xi^{-m-m'+\abs\beta}\abs{D^\alpha_zD^\beta_\xi(ab)}
    \le\sum_{\mu\le\alpha,\ \gamma\le\beta}\binom\alpha\mu\binom\beta\gamma
    \ang\xi^{-m+\abs\gamma}\abs{D^\mu_zD^\gamma_\xi a}\,
    \ang\xi^{-m'+\abs{\beta-\gamma}}\abs{D^{\alpha-\mu}_zD^{\beta-\gamma}_\xi b},
  \end{equation*}
  so $\norm{ab}_{N,m+m'}\le C_N\norm a_{N,m}\norm b_{N,m'}$.

  (iii) is immediate from the definition:
  $\norm{D^\alpha_za}_{N,m}\le\norm a_{N+\abs\alpha,m}$ and
  $\norm{D^\beta_\xi a}_{N,m-\abs\beta}\le\norm a_{N+\abs\beta,m}$.

  (iv) If $(a_k)$ is Cauchy in $S^m$ then, for each $\alpha,\beta$,
  $\ang\xi^{-m+\abs\beta}D^\alpha_zD^\beta_\xi a_k$ converges uniformly. Hence
  $a_k\to a$ in $\CI$ with all derivatives converging locally uniformly, and
  the uniform limits are $\ang\xi^{-m+\abs\beta}D^\alpha_zD^\beta_\xi a$; it
  follows that $a\in S^m$ and $a_k\to a$ in $S^m$.

  (v) $\ang\xi^{N}\abs{D^\alpha_zD^\beta_\xi a}\le\ang R^N\sup\abs{D^\alpha_zD^\beta_\xi a}$.
\end{proof}

In particular $S^m\cdot S^{m'}=S^{m+m'}$, since
$a=\ang\xi^{m}\cdot(\ang\xi^{-m}a)$. The inclusion
$S^m\hookrightarrow S^{m'}$, $m'>m$, is never dense. Instead we have the
following technical but very useful result.

\begin{lemma}\label{euclidean:lem:density}
  Let $\chi\in\CIc(\R^n)$ with $0\le\chi\le1$, $\chi(\xi)=1$ for
  $\abs\xi\le1$ and $\chi(\xi)=0$ for $\abs\xi\ge2$. For
  $a\in S^m(\R^p;\R^n)$ and $0<\eps\le1$ set
  $a_\eps(z,\xi)=\chi(\eps\xi)a(z,\xi)$. Then $a_\eps\in S^{-\infty}(\R^p;\R^n)$,
  the family $(a_\eps)$ is bounded in $S^m(\R^p;\R^n)$, and $a_\eps\to a$ in
  $S^{m'}(\R^p;\R^n)$ as $\eps\to0$ for every $m'>m$. In particular
  $S^{-\infty}(\R^p;\R^n)$ is dense in $S^m(\R^p;\R^n)$ in the topology of
  $S^{m'}(\R^p;\R^n)$ for every $m'>m$.
\end{lemma}
\begin{proof}
  The function $\chi(\eps\xi)$ is bounded in $S^0(\R^n)$ uniformly in
  $\eps\in(0,1]$: for $\beta\ne0$,
  $D^\beta_\xi(\chi(\eps\xi))=\eps^{\abs\beta}(D^\beta\chi)(\eps\xi)$ is
  supported in $1/\eps\le\abs\xi\le2/\eps$, where
  $\ang\xi\le1+\abs\xi\le3/\eps$, so it is bounded by
  $3^{\abs\beta}\sup\abs{D^\beta\chi}\,\ang\xi^{-\abs\beta}$. By
  \cref{euclidean:lem:symbol-basic}(ii) the family $a_\eps$ is bounded in
  $S^m$, and $a_\eps\in S^{-\infty}$ by \cref{euclidean:lem:symbol-basic}(v).
  Now $D^\alpha_zD^\beta_\xi(a-a_\eps)$ is bounded by
  $C_{\alpha\beta}\ang\xi^{m-\abs\beta}$ uniformly in $\eps$ and vanishes for
  $\abs\xi<1/\eps$. Hence, for $m'>m$,
  \begin{equation*}
    \sup\ang\xi^{-m'+\abs\beta}\abs{D^\alpha_zD^\beta_\xi(a-a_\eps)}
    \le C_{\alpha\beta}\sup_{\abs\xi\ge1/\eps}\ang\xi^{m-m'}
    \le C_{\alpha\beta}\,\eps^{m'-m}\to0.\qedhere
  \end{equation*}
\end{proof}

One cannot take $m'=m$ here, for essentially the same reason that $\cS(\R^n)$
is not dense in $\CIb(\R^n)$: a limit in $S^m$ of residual symbols satisfies
$\ang\xi^{-m}a\to0$ as $\abs\xi\to\infty$, so for instance
$1\in S^0(\R^p;\R^n)$ is not in the closure of $S^{-\infty}$ if $n>0$.

\subsection{Ellipticity of symbols}

\begin{lemma}\label{euclidean:lem:powers}
  Let $a\in S^m(\R^p;\R^n)$ and let $U\subset\R^p\times\R^n$ be an open set
  on which $\abs{a(z,\xi)}\ge\eps\ang\xi^m$ for some $\eps>0$. Then on $U$
  \begin{equation*}
    \abs{D^\alpha_zD^\beta_\xi(a^{-1})}\le C_{\alpha\beta}\ang\xi^{-m-\abs\beta},
  \end{equation*}
  and if moreover $a>0$ on $U$ then, for every $s\in\R$,
  $\abs{D^\alpha_zD^\beta_\xi(a^{s})}\le C_{\alpha\beta s}\ang\xi^{sm-\abs\beta}$
  on $U$. The constants depend only on $\eps$, $s$ and the norms of $a$. In
  particular, if $U=\R^p\times\R^n$, then $a^{-1}\in S^{-m}(\R^p;\R^n)$ and,
  when $a>0$, $a^s\in S^{sm}(\R^p;\R^n)$.
\end{lemma}
\begin{proof}
  Let $s\in\R$ (with $s=-1$ if $a$ is not positive). We claim that on $U$
  \begin{equation*}
    D^\alpha_zD^\beta_\xi(a^s)=\sum_{j=0}^{\abs\alpha+\abs\beta}a^{s-j}G_{j},
    \qquad G_j=G_{j}^{\alpha\beta}\in S^{jm-\abs\beta}(\R^p;\R^n),
  \end{equation*}
  where the $G_j$ are universal polynomials in the derivatives of $a$ (hence
  defined on all of $\R^p\times\R^n$). This holds for $\alpha=\beta=0$ with
  $G_0=1$. If it holds for $(\alpha,\beta)$ then
  \begin{equation*}
    D_{z_l}(a^{s-j}G_j)=a^{s-j-1}\bigl((s-j)(D_{z_l}a)G_j\bigr)+a^{s-j}D_{z_l}G_j,
  \end{equation*}
  with $(D_{z_l}a)G_j\in S^{(j+1)m-\abs\beta}$ and
  $D_{z_l}G_j\in S^{jm-\abs\beta}$; similarly for $D_{\xi_l}$, where
  $(D_{\xi_l}a)G_j\in S^{(j+1)m-\abs\beta-1}$ and
  $D_{\xi_l}G_j\in S^{jm-\abs\beta-1}$. This proves the claim. On $U$,
  $\abs a^{s-j}\le C\ang\xi^{(s-j)m}$: for $s-j\ge0$ by the upper bound
  $\abs a\le C\ang\xi^m$, for $s-j<0$ by the lower bound. Hence
  $\abs{a^{s-j}G_j}\le C\ang\xi^{(s-j)m+jm-\abs\beta}=C\ang\xi^{sm-\abs\beta}$.
\end{proof}

\begin{definition}\label{euclidean:def:elliptic-symbol}
  A symbol $a\in S^m(\R^p;\R^n)$ is \emph{elliptic} if there are $\eps>0$ and
  $C$ such that
  \begin{equation}\label{euclidean:eq:elliptic-symbol}
    \abs{a(z,\xi)}\ge\eps\ang\xi^m\quad\text{for all }z\in\R^p,\ \abs\xi\ge C.
  \end{equation}
\end{definition}

Note that $\eps$ and $C$ are independent of $z$, so this is a notion of
\emph{uniform} ellipticity. Condition \eqref{euclidean:eq:elliptic-symbol} is
equivalent to the existence of $\eps>0$ and $C'$ with
$\abs{a(z,\xi)}\ge\eps\ang\xi^m-C'\ang\xi^{m-1}$ everywhere. Indeed, the
latter gives $\abs a\ge\frac\eps2\ang\xi^m$ once $\ang\xi\ge2C'/\eps$;
conversely, if \eqref{euclidean:eq:elliptic-symbol} holds then for
$\abs\xi\le C$ we have $\eps\ang\xi^m\le\eps\ang C\ang\xi^{m-1}$, so
$\abs a\ge\eps\ang\xi^m-\eps\ang C\ang\xi^{m-1}$ everywhere.

\begin{lemma}\label{euclidean:lem:elliptic-symbol}
  If $a\in S^m(\R^p;\R^n)$ is elliptic there exists
  $b\in S^{-m}(\R^p;\R^n)$ such that $ab-1\in S^{-\infty}(\R^p;\R^n)$.
  Conversely, if $ab-1\in S^{-\infty}$ for some $b\in S^{-m}$, then $a$ is
  elliptic.
\end{lemma}
\begin{proof}
  Let $C$ be as in \eqref{euclidean:eq:elliptic-symbol} and choose
  $\psi\in\CI(\R^n)$ with $\psi(\xi)=0$ for $\abs\xi\le C+1$ and
  $\psi(\xi)=1$ for $\abs\xi\ge C+2$; then $\psi\in S^0(\R^n)$. Set
  $b=\psi/a$ where $\psi\ne0$ and $b=0$ elsewhere; $b$ is smooth since it
  vanishes on a neighbourhood of $\{\abs\xi\le C+1\}$ and $a\neq 0$ on
  $\abs\xi>C$. Applying \cref{euclidean:lem:powers} on
  $U=\{\abs\xi>C\}$ and Leibniz' formula shows $b\in S^{-m}$, and
  $ab-1=\psi-1$ is smooth, independent of $z$ and supported in
  $\abs\xi\le C+2$, hence residual.

  Conversely, if $ab=1+e$ with $e\in S^{-\infty}$, choose $C$ with
  $\abs e\le\frac12$ for $\abs\xi\ge C$. Since $\abs b\le C'\ang\xi^{-m}$, for
  $\abs\xi\ge C$ we get
  $\frac12\le\abs a\abs b\le C'\ang\xi^{-m}\abs a$, which is
  \eqref{euclidean:eq:elliptic-symbol} with $\eps=1/(2C')$.
\end{proof}

\subsection{Asymptotic summation}

Suppose $a_j\in S^{m-j}(\R^p;\R^n)$, $j\in\N_0$. The decreasing orders mean
that these symbols become smaller and smaller as $\abs\xi\to\infty$, but the
series $\sum_ja_j$ need not converge. Nevertheless it can be ``summed''
asymptotically.

\begin{definition}\label{euclidean:def:asymptotic-sum}
  Let $m_j\in\R$ be a sequence decreasing to $-\infty$ and
  $a_j\in S^{m_j}(\R^p;\R^n)$. We write
  \begin{equation*}
    a\sim\sum_{j=0}^\infty a_j
  \end{equation*}
  if $a\in S^{m_0}(\R^p;\R^n)$ and
  $a-\sum_{j<N}a_j\in S^{m_N}(\R^p;\R^n)$ for every $N$; $a$ is then called
  an \emph{asymptotic sum} of the series. Most often $m_j=m-j$.
\end{definition}

\begin{proposition}[Borel's lemma for symbols]\label{euclidean:prop:borel}
  Every series $\sum_ja_j$ with $a_j\in S^{m_j}(\R^p;\R^n)$, $m_j\downarrow-\infty$,
  has an asymptotic sum, and the asymptotic sum is unique up to an additive
  term in $S^{-\infty}(\R^p;\R^n)$.
\end{proposition}
\begin{proof}
  Uniqueness: if $a$ and $a'$ are both asymptotic sums then
  $a-a'=(a-\sum_{j<N}a_j)-(a'-\sum_{j<N}a_j)\in S^{m_N}$ for every $N$, so
  $a-a'\in S^{-\infty}$.

  Existence: we cut off each term near the origin in $\xi$, more and more
  drastically as $j$ increases. Choose $\phi\in\CI(\R^n)$ with
  $\phi(\xi)=0$ for $\abs\xi\le1$ and $\phi(\xi)=1$ for $\abs\xi\ge2$
  (e.g.\ $\phi=1-\chi$ with $\chi$ as in \cref{euclidean:lem:density}).
  As in the proof of \cref{euclidean:lem:density}, $\phi(\eps\xi)$ is
  bounded in $S^0(\R^n)$ for $\eps\in(0,1]$, so
  $\norm{\phi(\eps\xi)a_j}_{N,m_j}\le C_{N,j}$ with $C_{N,j}$ independent
  of $\eps$. Moreover $\phi(\eps\xi)a_j$ vanishes for $\abs\xi<1/\eps$, where
  $\ang\xi^{-1}\le\eps$; replacing $m_j$ by $m_j+1$ in the norm therefore
  gains a factor $\eps$:
  \begin{equation}\label{euclidean:eq:borel-gain}
    \norm{\phi(\eps\xi)a_j}_{N,m_j+1}\le C_{N,j}\,\eps.
  \end{equation}
  Choose $\eps_j\in(0,1]$ with
  $C_{N,j}\eps_j\le2^{-j}$ for all $N\le j$, and set
  \begin{equation*}
    a(z,\xi)=\sum_{j=0}^\infty\phi(\eps_j\xi)a_j(z,\xi).
  \end{equation*}
  The $\eps_j$ may be taken decreasing to $0$, so on each set
  $\{\abs\xi\le R\}$ only finitely many terms are non-zero and $a$ is a
  well-defined smooth function. Fix $M\in\N_0$. Since $m_j\downarrow-\infty$
  there is $J_M$ with $m_j+1\le m_M$ for $j\ge J_M$; for such $j$ and
  $j\ge N$,
  $\norm{\phi(\eps_j\xi)a_j}_{N,m_M}\le\norm{\phi(\eps_j\xi)a_j}_{N,m_j+1}\le2^{-j}$.
  Hence $\sum_{j\ge M}\phi(\eps_j\xi)a_j$ converges absolutely in each norm
  of the Fr\'echet space $S^{m_M}$ and so converges there; its limit is the
  pointwise sum. Thus
  \begin{equation*}
    a-\sum_{j<M}a_j=\sum_{j<M}(\phi(\eps_j\xi)-1)a_j+\sum_{j\ge M}\phi(\eps_j\xi)a_j
    \in S^{-\infty}+S^{m_M},
  \end{equation*}
  where the first sum is residual by \cref{euclidean:lem:symbol-basic}(v).
\end{proof}

The following simple observations allow asymptotic sums to be manipulated
formally.

\begin{lemma}\label{euclidean:lem:asymptotic-manip}
  \begin{enumerate}[label=\textup{(\roman*)}]
  \item If $T$ is a linear map with $T(S^{k}(\R^p;\R^n))\subset S^{k+\mu}(\R^{p'};\R^{n})$ for all
    $k$ \textup{(}e.g.\ multiplication by an element of $S^\mu$,
    differentiation, complex conjugation, restriction to a linear subspace of
    $\R^p$\textup{)} and $a\sim\sum_ja_j$, then $Ta\sim\sum_jTa_j$.
  \item If $a\sim\sum_ja_j$ and $b\sim\sum_kb_k$, with $a_j\in S^{m-j}$,
    $b_k\in S^{m'-k}$, then $ab\sim\sum_l\sum_{j+k=l}a_jb_k$.
  \item \textup{(Rearrangement)} Let $a_{j,k}\in S^{m-j-k}(\R^p;\R^n)$ for
    $j,k\in\N_0$. If $c_j\sim\sum_ka_{j,k}$ for each $j$ and
    $c\sim\sum_jc_j$, then $c\sim\sum_l\bigl(\sum_{j+k=l}a_{j,k}\bigr)$.
  \end{enumerate}
\end{lemma}
\begin{proof}
  (i) $Ta-\sum_{j<N}Ta_j=T(a-\sum_{j<N}a_j)\in S^{m_N+\mu}$.
  (ii) $ab-\sum_{j+k<N}a_jb_k=(a-\sum_{j<N}a_j)b+\sum_{j<N}a_j(b-\sum_{k<N-j}b_k)
  \in S^{m+m'-N}$.
  (iii) $c-\sum_{j+k<N}a_{j,k}=(c-\sum_{j<N}c_j)+\sum_{j<N}(c_j-\sum_{k<N-j}a_{j,k})
  \in S^{m-N}$.
\end{proof}

% ======================================================================
\section{Amplitudes and pseudodifferential operators}\label{euclidean:sec:operators}

We now interpret \eqref{euclidean:eq:amplitude-operator}. It turns out to be
convenient to allow the amplitude $a(x,y,\xi)$ to grow polynomially in the
direction of $x-y$; this extra freedom makes no difference to the class of
operators (\cref{euclidean:thm:reduction}) but is very useful, for instance
for weighted estimates.

\begin{definition}\label{euclidean:def:amplitude}
  For $m,w\in\R$ let $\Amp mw(\R^n)$ be the space of
  $a\in\CI(\R^n_x\times\R^n_y\times\R^n_\xi)$ such that
  \begin{equation}\label{euclidean:eq:amplitude-est}
    \abs{D^\alpha_xD^\beta_yD^\gamma_\xi a(x,y,\xi)}\le
    C_{\alpha\beta\gamma}\ang{x-y}^w\ang\xi^{m-\abs\gamma}
    \quad\text{for all }\alpha,\beta,\gamma\in\N_0^n,
  \end{equation}
  with the norms $\norm a_{N,m,w}$, the best constants in
  \eqref{euclidean:eq:amplitude-est} for $\abs\alpha+\abs\beta+\abs\gamma\le N$.
  Elements of $\Amp mw(\R^n)$ are called \emph{amplitudes} of order $m$ and
  weight $w$. We set $\Amp{-\infty}{w}(\R^n)=\bigcap_m\Amp mw(\R^n)$.
\end{definition}

\begin{lemma}\label{euclidean:lem:amplitude-char}
  Multiplication by $\ang{x-y}^w$ is an isomorphism of Fr\'echet spaces
  $S^m(\R^{2n}_{(x,y)};\R^n_\xi)\to\Amp mw(\R^n)$. More generally, call
  $f\in\CI(\R^{2n})$ \emph{of weight $w'$} if
  $\abs{D^\alpha_xD^\beta_yf}\le C_{\alpha\beta}\ang{x-y}^{w'}$ for all
  $\alpha,\beta$; then multiplication by such $f$ maps $\Amp mw(\R^n)$
  continuously into $\Amp m{w+w'}(\R^n)$.
\end{lemma}
\begin{proof}
  The second statement is immediate from Leibniz' formula. The function
  $\ang z^w$ is a symbol of order $w$ in $z\in\R^n$ (with no base variables)
  by \cref{euclidean:lem:symbol-basic}(i), so
  $\abs{D^\alpha_z\ang z^{w}}\le C\ang z^{w-\abs\alpha}\le C\ang z^w$. Hence
  $\ang{x-y}^{\pm w}$ is of weight $\pm w$, and multiplication by
  $\ang{x-y}^{w}$ and by $\ang{x-y}^{-w}$ are mutually inverse continuous
  maps between $S^m(\R^{2n};\R^n)=\Amp m0(\R^n)$ and $\Amp mw(\R^n)$.
\end{proof}

Thus $\Amp mw(\R^n)=\ang{x-y}^wS^m(\R^{2n};\R^n)$, and the results of
\cref{euclidean:sec:symbols} transfer to amplitudes. In particular, by
\cref{euclidean:lem:density}, if $a\in\Amp mw(\R^n)$ and $\chi$ is as there,
then $\chi(\eps\xi)a\in\Amp{-\infty}w(\R^n)$ is bounded in $\Amp mw(\R^n)$ and
converges to $a$ in $\Amp{m'}w(\R^n)$ as $\eps\to0$, for every $m'>m$. Note
also the following elementary facts:
\begin{itemize}
\item $\Amp mw\subset\Amp{m'}{w'}$ if $m\le m'$ and $w\le w'$;
  $\Amp mw\cdot\Amp{m'}{w'}\subset\Amp{m+m'}{w+w'}$;
\item $D_{x_j}$, $D_{y_j}$ preserve $\Amp mw$, $D_{\xi_j}$ maps it to
  $\Amp{m-1}w$ and multiplication by $\xi_j$ maps it to $\Amp{m+1}w$;
\item $(x-y)^\alpha$ is of weight $\abs\alpha$, so
  $(x-y)^\alpha\Amp mw\subset\Amp m{w+\abs\alpha}$;
\item if $a\in S^m(\R^n;\R^n)$ then both $a(x,\xi)$ and $a(y,\xi)$ lie in
  $\Amp m0(\R^n)$; and if $a\in\Amp mw$ then
  $a(x,x,\xi)\in S^m(\R^n;\R^n)$.
\end{itemize}

\subsection{The kernel as an oscillatory integral}

If $a\in\Amp mw(\R^n)$ with $m<-n$ then
\begin{equation*}
  K(x,y)=\int e^{i(x-y)\cdot\xi}a(x,y,\xi)\,\dbar\xi
\end{equation*}
converges absolutely and locally uniformly, so defines a continuous function
with $\abs{K(x,y)}\le C\ang{x-y}^w$. For general $m$ we integrate by parts.
Let $\Delta_x=D_{x_1}^2+\dots+D_{x_n}^2$ be the (positive) Laplacian. Then
\begin{equation}\label{euclidean:eq:laplace-exp}
  (1+\Delta_x)^qe^{i(x-y)\cdot\xi}=\ang\xi^{2q}e^{i(x-y)\cdot\xi},\qquad
  (1+\Delta_y)^qe^{i(x-y)\cdot\xi}=\ang\xi^{2q}e^{i(x-y)\cdot\xi}.
\end{equation}

\begin{proposition}\label{euclidean:prop:osc-int}
  Let $a\in\Amp mw(\R^n)$ and let $q\in\N_0$ satisfy $m-2q<-n$. For
  $\phi\in\cS(\R^{2n})$ the integral
  \begin{equation}\label{euclidean:eq:I-def}
    I(a)(\phi)=\int_{\R^{3n}}e^{i(x-y)\cdot\xi}\ang\xi^{-2q}(1+\Delta_x)^q
    \bigl[a(x,y,\xi)\phi(x,y)\bigr]\,dx\,dy\,\dbar\xi
  \end{equation}
  converges absolutely, is independent of the choice of such $q$, and
  satisfies
  \begin{equation}\label{euclidean:eq:I-est}
    \abs{I(a)(\phi)}\le C\norm a_{2q,m,w}\norm\phi_k
  \end{equation}
  for some $k=k(q,w,n)$ and $C=C(q,w,n,m)$. Thus $I(a)\in\cSp(\R^{2n})$ and
  $a\mapsto I(a)$ is continuous in the sense of \eqref{euclidean:eq:I-est}.
  Moreover:
  \begin{enumerate}[label=\textup{(\roman*)}]
  \item if $m<-n$, $I(a)$ is the continuous function
    $K(x,y)=\int e^{i(x-y)\cdot\xi}a(x,y,\xi)\,\dbar\xi$;
  \item if $a_k$ is bounded in $\Amp mw(\R^n)$ and $a_k\to a$ in
    $\Amp{m'}w(\R^n)$ for some $m'>m$, then $I(a_k)(\phi)\to I(a)(\phi)$ for
    every $\phi\in\cS(\R^{2n})$; in particular, for $\chi$ as in
    \cref{euclidean:lem:density},
    \begin{equation*}
      I(a)(\phi)=\lim_{\eps\to0}\int e^{i(x-y)\cdot\xi}\chi(\eps\xi)a(x,y,\xi)
      \phi(x,y)\,dx\,dy\,\dbar\xi;
    \end{equation*}
  \item the same value is obtained if $(1+\Delta_x)^q$ is replaced by
    $(1+\Delta_y)^q$ in \eqref{euclidean:eq:I-def}.
  \end{enumerate}
\end{proposition}

Part (ii) justifies calling $I(a)$ an \emph{oscillatory integral}: it is
the limit of the absolutely convergent integrals obtained by cutting off the
amplitude at large frequencies, and it is the unique extension of the
absolutely convergent case that is continuous in the sense of (ii).

\begin{proof}
  By Leibniz' formula $(1+\Delta_x)^q[a\phi]$ is a linear combination of
  terms $D^\mu_xa\,D^\nu_x\phi$ with $\abs\mu+\abs\nu\le2q$, and
  $\abs{D^\mu_xa}\le\norm a_{2q,m,w}\ang{x-y}^w\ang\xi^m$. Since
  $\ang{x-y}\le\sqrt2\ang{(x,y)}$ we have
  $\ang{x-y}^w\le2^{\abs w/2}\ang{(x,y)}^{\abs w}$. Hence the integrand in
  \eqref{euclidean:eq:I-def} is bounded by
  \begin{equation*}
    C\norm a_{2q,m,w}\ang\xi^{m-2q}\ang{(x,y)}^{\abs w}
    \sum_{\abs\nu\le2q}\abs{D^\nu\phi(x,y)}
    \le C'\norm a_{2q,m,w}\norm\phi_k\ang\xi^{m-2q}\ang{(x,y)}^{-2n-1}
  \end{equation*}
  for $k\ge2q+\abs w+2n+1$, which is integrable since $m-2q<-n$. This proves
  absolute convergence and \eqref{euclidean:eq:I-est}.

  If $q$ and $q+1$ are both admissible, put
  $F=\ang\xi^{-2q-2}(1+\Delta_x)^q[a\phi]$. For fixed $(y,\xi)$, $F$ and all
  its $x$-derivatives are rapidly decreasing in $x$, so integration by
  parts and \eqref{euclidean:eq:laplace-exp} give
  \begin{equation*}
    \int e^{i(x-y)\cdot\xi}(1+\Delta_x)F\,dx=\int\bigl((1+\Delta_x)e^{i(x-y)\cdot\xi}\bigr)F\,dx
    =\ang\xi^2\int e^{i(x-y)\cdot\xi}F\,dx.
  \end{equation*}
  Integrating in $(y,\xi)$ (both sides being absolutely integrable) shows that
  the integrals for $q+1$ and $q$ agree. So $I(a)$ is independent of $q$.

  (i) If $m<-n$ we may take $q=0$, and then Fubini's theorem gives
  $I(a)(\phi)=\int K\phi$.

  (ii) Choose $q$ with $m'-2q<-n$ and apply \eqref{euclidean:eq:I-est} with
  $m'$ in place of $m$ to $a-a_k$. The particular case follows from the
  remarks after \cref{euclidean:lem:amplitude-char} and (i).

  (iii) Denote by $I'(a)(\phi)$ the integral with $\Delta_y$ in place of
  $\Delta_x$. Exactly the same arguments show that it converges absolutely,
  is independent of $q$ and satisfies \eqref{euclidean:eq:I-est}, hence (ii).
  For $a\in\Amp{-\infty}w$ both $I(a)(\phi)$ and $I'(a)(\phi)$ equal
  $\int e^{i(x-y)\cdot\xi}a\phi\,dx\,dy\,\dbar\xi$ (take $q=0$). By (ii), applied
  to $\chi(\eps\xi)a$, they agree for all $a$.
\end{proof}

Distributional operations on $I(a)$ correspond to simple operations on the
amplitude.

\begin{lemma}\label{euclidean:lem:kernel-identities}
  For $a\in\Amp mw(\R^n)$ and $j=1,\dots,n$, the following identities hold in
  $\cSp(\R^{2n})$:
  \begin{gather*}
    D_{x_j}I(a)=I\bigl((\xi_j+D_{x_j})a\bigr),\qquad
    D_{y_j}I(a)=I\bigl((D_{y_j}-\xi_j)a\bigr),\\
    (x_j-y_j)I(a)=I\bigl((x_j-y_j)a\bigr)=I(-D_{\xi_j}a).
  \end{gather*}
  Moreover $fI(a)=I(fa)$ for every $f$ of weight $w'$ in the sense of
  \cref{euclidean:lem:amplitude-char}. In particular, for every
  $\alpha\in\N_0^n$, $(x-y)^\alpha I(a)=I\bigl((-D_\xi)^\alpha a\bigr)$.
\end{lemma}
\begin{proof}
  Such an $f$ and all its derivatives are polynomially bounded, so
  $f\phi\in\cS(\R^{2n})$ for $\phi\in\cS(\R^{2n})$, and the integrals defining
  $I(a)(f\phi)$ and $I(fa)(\phi)$ by \eqref{euclidean:eq:I-def} (with the same
  $q$) are identical. This gives $fI(a)=I(fa)$, and in particular the first
  equality in the second line.

  For the remaining identities, suppose first that $a\in\Amp{-\infty}w$. Then
  all integrals below converge absolutely and
  \begin{equation*}
    D_{x_j}I(a)(\phi)=-\int e^{i(x-y)\cdot\xi}a\,D_{x_j}\phi
    =\int D_{x_j}\bigl(e^{i(x-y)\cdot\xi}a\bigr)\phi
    =\int e^{i(x-y)\cdot\xi}(\xi_ja+D_{x_j}a)\phi,
  \end{equation*}
  integrating by parts in $x$. The identity for $D_{y_j}$ is the same.
  Since $(x_j-y_j)e^{i(x-y)\cdot\xi}=D_{\xi_j}e^{i(x-y)\cdot\xi}$, integration
  by parts in $\xi$ (the boundary terms vanish because $a$ is rapidly
  decreasing in $\xi$) gives $I((x_j-y_j)a)=I(-D_{\xi_j}a)$.
  For general $a$, apply these identities to $\chi(\eps\xi)a$ and let
  $\eps\to0$, using \cref{euclidean:prop:osc-int}(ii) for both sides: the
  amplitudes involved are bounded in $\Amp{m+1}{w+1}$ and converge in
  $\Amp{m'+1}{w+1}$ for $m'>m$. The last statement follows by induction.
\end{proof}

\subsection{Pseudodifferential operators}

\begin{proposition}\label{euclidean:prop:action-S}
  Let $a\in\Amp mw(\R^n)$. For $u\in\cS(\R^n)$ define $\Op(a)u\in\cSp(\R^n)$
  by
  \begin{equation}\label{euclidean:eq:Op-def}
    \pair{\Op(a)u}v=I(a)(v\otimes u),\qquad v\in\cS(\R^n),
  \end{equation}
  where $(v\otimes u)(x,y)=v(x)u(y)$. Then $\Op(a)u\in\cS(\R^n)$ and
  $\Op(a):\cS(\R^n)\to\cS(\R^n)$ is continuous. For $q\in\N_0$ with
  $m-2q<-n$,
  \begin{equation}\label{euclidean:eq:Op-formula}
    \Op(a)u(x)=\int e^{i(x-y)\cdot\xi}\ang\xi^{-2q}(1+\Delta_y)^q
    \bigl[a(x,y,\xi)u(y)\bigr]\,dy\,\dbar\xi,
  \end{equation}
  an absolutely convergent integral; if $m<-n$ this is just
  $\int e^{i(x-y)\cdot\xi}a(x,y,\xi)u(y)\,dy\,\dbar\xi$. In general
  \begin{equation*}
    \Op(a)u(x)=\lim_{\eps\to0}\int e^{i(x-y)\cdot\xi}\chi(\eps\xi)a(x,y,\xi)u(y)\,dy\,\dbar\xi.
  \end{equation*}
\end{proposition}
\begin{proof}
  Let $F(x)$ denote the right side of \eqref{euclidean:eq:Op-formula}. As in
  the proof of \cref{euclidean:prop:osc-int}, but now using Peetre's
  inequality $\ang{x-y}^w\le2^{\abs w/2}\ang x^{\abs w}\ang y^{\abs w}$, its
  integrand is bounded by
  $C\norm a_{2q,m,w}\ang x^{\abs w}\ang\xi^{m-2q}\ang y^{\abs w}
  \sum_{\abs\nu\le2q}\abs{D^\nu u(y)}$. Hence the integral converges
  absolutely, $F$ is continuous by dominated convergence, and
  \begin{equation}\label{euclidean:eq:Op-bound}
    \abs{F(x)}\le C\norm a_{2q,m,w}\norm u_k\ang x^{\abs w}
  \end{equation}
  for some $k$. By \cref{euclidean:prop:osc-int}(iii) and Fubini's theorem,
  $I(a)(v\otimes u)=\int v(x)F(x)\,dx$, so $\Op(a)u=F$. The limit formula
  follows in the same way from \cref{euclidean:prop:osc-int}(ii), the
  convergence being locally uniform in $x$ by dominated convergence.

  To see that $\Op(a)u\in\cS$, note that \eqref{euclidean:eq:Op-def} and
  \cref{euclidean:lem:kernel-identities} give
  \begin{equation*}
    D_{x_j}\Op(a)u=\Op\bigl((\xi_j+D_{x_j})a\bigr)u,\qquad
    x_j\Op(a)u=\Op(a)(x_ju)+\Op(-D_{\xi_j}a)u.
  \end{equation*}
  For instance
  $\pair{D_{x_j}\Op(a)u}v=-I(a)(D_{x_j}v\otimes u)=(D_{x_j}I(a))(v\otimes u)$,
  and $(x_j-y_j)(v\otimes u)=(x_jv)\otimes u-v\otimes(x_ju)$. By induction,
  for all $\alpha,\beta$, $x^\alpha D^\beta\Op(a)u$ is a finite sum of terms
  $\Op(a')(x^{\alpha'}u)$ with $a'\in\Amp{m+\abs\beta}w$, $\abs{\alpha'}\le\abs\alpha$,
  the $a'$ depending continuously on $a$. By \eqref{euclidean:eq:Op-bound}
  each such term is a continuous function bounded by
  $C\norm a_{N,m,w}\norm u_{k'}\ang x^{\abs w}$. Since all distributional
  derivatives of $\Op(a)u$ are continuous, $\Op(a)u\in\CI(\R^n)$; and since
  $\ang x^{2L}$ is a polynomial in $x$, we get
  $\sup\ang x^{2L-\abs w}\abs{D^\beta\Op(a)u}\le C\norm a_{N,m,w}\norm u_{k''}$
  for all $L$ and $\beta$. Thus $\Op(a)u\in\cS(\R^n)$ and $\Op(a)$ is
  continuous on $\cS(\R^n)$.
\end{proof}

\begin{definition}\label{euclidean:def:psido}
  For $m\in\R$, $\Psido m(\R^n)$ is the space of operators
  $\Op(a):\cS(\R^n)\to\cS(\R^n)$ with $a\in\Amp mw(\R^n)$ for some $w\in\R$.
  Its elements are the \emph{pseudodifferential operators of order $m$} (of
  type $(1,0)$, with uniformly bounded coefficients). We set
  $\Psido\infty(\R^n)=\bigcup_m\Psido m(\R^n)$ and
  $\Psido{-\infty}(\R^n)=\bigcap_m\Psido m(\R^n)$. The distribution
  $K_A=I(a)$ is the \emph{Schwartz kernel} of $A=\Op(a)$.
\end{definition}

Since $\Amp mw\subset\Amp{m'}{w'}$ for $m\le m'$, $w\le w'$, each
$\Psido m(\R^n)$ is a linear space and $\Psido m(\R^n)\subset\Psido{m'}(\R^n)$
for $m\le m'$. By \eqref{euclidean:eq:Op-def}, $K_A$ is determined by $A$ on
the products $v\otimes u$, and these span a dense subspace of $\cS(\R^{2n})$
(see \cref{ch:distributions}), so $K_A$ is determined by $A$ even though the
amplitude is not.

The two extreme cases are amplitudes independent of $y$ and amplitudes
independent of $x$. For $a,b\in S^m(\R^n;\R^n)$ we write
\begin{equation*}
  \Op(a)=\Op\bigl(a(x,\xi)\bigr),\qquad\Op_R(b)=\Op\bigl(b(y,\xi)\bigr),
\end{equation*}
the \emph{left} and \emph{right quantizations}; they lie in $\Psido m(\R^n)$.

\begin{proposition}\label{euclidean:prop:left-right}
  Let $a,b\in S^m(\R^n;\R^n)$ and $u\in\cS(\R^n)$. Then
  \begin{equation}\label{euclidean:eq:left-quant}
    \Op(a)u(x)=\int e^{ix\cdot\xi}a(x,\xi)\hat u(\xi)\,\dbar\xi,
  \end{equation}
  \begin{equation}\label{euclidean:eq:right-quant}
    \widehat{\Op_R(b)u}(\xi)=\int e^{-iy\cdot\xi}b(y,\xi)u(y)\,dy,
  \end{equation}
  Moreover, if $a\in S^m(\R^n;\R^n)$ and $b\in S^{m'}(\R^n;\R^n)$ then
  $a(x,\xi)b(y,\xi)\in\Amp{m+m'}0(\R^n)$ and
  $\Op(a)\Op_R(b)=\Op\bigl(a(x,\xi)b(y,\xi)\bigr)\in\Psido{m+m'}(\R^n)$.
\end{proposition}
\begin{proof}
  If $a\in S^{-\infty}$, \eqref{euclidean:eq:left-quant} follows from the
  absolutely convergent formula of \cref{euclidean:prop:action-S} and Fubini.
  In general, let $a_\eps=\chi(\eps\xi)a$. Then $\Op(a_\eps)u(x)\to\Op(a)u(x)$
  by \cref{euclidean:prop:action-S}, while
  $\int e^{ix\cdot\xi}a_\eps\hat u\,\dbar\xi\to\int e^{ix\cdot\xi}a\hat u\,\dbar\xi$
  by dominated convergence, since $\abs{a_\eps}\le C\ang\xi^m$ uniformly
  and $\hat u\in\cS$.

  For \eqref{euclidean:eq:right-quant} let $g(\xi)$ be the right side and
  $B=\Op_R(b)$. If $b\in S^{-\infty}$, $g$ is integrable and Fubini gives
  $Bu(x)=\int e^{ix\cdot\xi}g(\xi)\,\dbar\xi$, so $\widehat{Bu}=g$ by Fourier
  inversion. In general, with $b_\eps=\chi(\eps\xi)b$ and $B_\eps=\Op_R(b_\eps)$,
  we have $\widehat{B_\eps u}=g_\eps\to g$ pointwise by dominated convergence;
  on the other hand $B_\eps u\to Bu$ in $\cSp(\R^n)$, by
  \cref{euclidean:prop:osc-int}(ii), so $\widehat{B_\eps u}\to\widehat{Bu}$
  in $\cSp$. Since $g_\eps$ is bounded by $C\ang\xi^{m}\norm u_{L^1}$
  uniformly in $\eps$, also $g_\eps\to g$ in $\cSp$, and so $\widehat{Bu}=g$.

  Finally let $c(x,y,\xi)=a(x,\xi)b(y,\xi)\in\Amp{m+m'}0$ and choose $q$ with
  $m+m'-2q<-n$. By \eqref{euclidean:eq:Op-formula} and Fubini,
  \begin{equation*}
    \Op(c)u(x)=\int e^{ix\cdot\xi}a(x,\xi)\ang\xi^{-2q}
    \Bigl[\int e^{-iy\cdot\xi}(1+\Delta_y)^q\bigl(b(y,\xi)u(y)\bigr)dy\Bigr]\dbar\xi .
  \end{equation*}
  Integrating by parts in $y$ (for fixed $\xi$, $bu$ is rapidly decreasing
  with all derivatives) the inner integral is $\ang\xi^{2q}g(\xi)$, so by
  \eqref{euclidean:eq:left-quant} and \eqref{euclidean:eq:right-quant},
  $\Op(c)u(x)=\int e^{ix\cdot\xi}a(x,\xi)\widehat{Bu}(\xi)\,\dbar\xi=\Op(a)Bu(x)$.
\end{proof}

\begin{example}\label{euclidean:ex:diff-ops}
  Let $P=\sum_{\abs\alpha\le m}p_\alpha(x)D^\alpha$ be a differential operator
  with coefficients $p_\alpha\in\CIb(\R^n)$; we denote the space of these by
  $\Diff m_{\mathrm b}(\R^n)$. Its characteristic polynomial
  $p(x,\xi)=\sum_\alpha p_\alpha(x)\xi^\alpha$ lies in $S^m(\R^n;\R^n)$ and,
  since $\widehat{D^\alpha u}=\xi^\alpha\hat u$, \eqref{euclidean:eq:left-quant}
  and Fourier inversion give $P=\Op(p)$. Thus
  $\Diff m_{\mathrm b}(\R^n)\subset\Psido m(\R^n)$. Similarly, multiplication
  by $f\in\CIb(\R^n)$ is $\Op(f)\in\Psido0(\R^n)$, and the Fourier multiplier
  $\ang D^s=\Op(\ang\xi^s)$, $\widehat{\ang D^su}=\ang\xi^s\hat u$, lies in
  $\Psido s(\R^n)$ for every $s\in\R$.
\end{example}

\subsection{Adjoints and the action on tempered distributions}

\begin{proposition}\label{euclidean:prop:adjoint}
  If $A=\Op(a)\in\Psido m(\R^n)$ with $a\in\Amp mw(\R^n)$, then
  $A^*=\Op(a^*)$ with $a^*(x,y,\xi)=\overline{a(y,x,\xi)}\in\Amp mw(\R^n)$
  satisfies
  \begin{equation*}
    \inner{Au}v=\inner u{A^*v}\quad\text{for all }u,v\in\cS(\R^n).
  \end{equation*}
  In particular $A^*\in\Psido m(\R^n)$, and the formal adjoint is
  determined by this identity. Consequently, the transpose
  $A^t\phi=\overline{A^*\bar\phi}$, which satisfies $\int(Au)\phi=\int u(A^t\phi)$
  for $u,\phi\in\cS$, is also continuous on $\cS(\R^n)$.
\end{proposition}
\begin{proof}
  Clearly $a^*\in\Amp mw$, depending continuously on $a$ in each $\Amp{m'}w$.
  We must show $I(a)(\bar v\otimes u)=\overline{I(a^*)(\bar u\otimes v)}$,
  since the left side is $\inner{Au}v$ and the right side is
  $\overline{\inner{A^*v}u}=\inner u{A^*v}$. For $a\in\Amp{-\infty}w$ this
  follows from the absolutely convergent expressions:
  \begin{equation*}
    \overline{I(a^*)(\bar u\otimes v)}
    =\int e^{-i(x-y)\cdot\xi}a(y,x,\xi)u(x)\overline{v(y)}\,dx\,dy\,\dbar\xi
    =\int e^{i(x-y)\cdot\xi}a(x,y,\xi)\overline{v(x)}u(y)\,dx\,dy\,\dbar\xi,
  \end{equation*}
  after exchanging the names of $x$ and $y$. The general case follows by
  replacing $a$ by $\chi(\eps\xi)a$ and using
  \cref{euclidean:prop:osc-int}(ii). Uniqueness of $A^*$ is clear since
  $\cS$ is dense in $L^2$.
\end{proof}

This shows one advantage of allowing general amplitudes: closure under
adjoints is immediate. It also allows every $A\in\Psido m(\R^n)$ to be
extended to tempered distributions by duality:
\begin{equation}\label{euclidean:eq:S'-action}
  (Au)(\phi)=u(A^t\phi),\qquad u\in\cSp(\R^n),\ \phi\in\cS(\R^n).
\end{equation}
Since $A^t$ is continuous on $\cS$, this defines $Au\in\cSp(\R^n)$, the map
$A:\cSp(\R^n)\to\cSp(\R^n)$ is continuous (for the weak topology), and it
extends $A$ on $\cS(\R^n)$; as $\cS(\R^n)$ is dense in $\cSp(\R^n)$ this
extension is unique. Since $(AB)^t=B^tA^t$, the extensions satisfy
$(AB)u=A(Bu)$ for $u\in\cSp$. For example
$\widehat{\ang D^su}=\ang\xi^s\hat u$ continues to hold for $u\in\cSp(\R^n)$.

% ======================================================================
\section{Reduction to the left symbol}\label{euclidean:sec:reduction}

The main step towards the composition theorem is to show that every
pseudodifferential operator can be written, uniquely, as the left
quantization of a symbol independent of $y$:

\begin{theorem}[Reduction]\label{euclidean:thm:reduction}
  For each $A\in\Psido m(\R^n)$ there is a unique
  $\sigma_L(A)\in S^m(\R^n;\R^n)$ with $A=\Op(\sigma_L(A))$. If
  $A=\Op(a)$ with $a\in\Amp mw(\R^n)$ then
  \begin{equation}\label{euclidean:eq:reduction}
    \sigma_L(A)(x,\xi)\sim\sum_{\alpha}\frac{i^{\abs\alpha}}{\alpha!}
    D^\alpha_\xi D^\alpha_ya(x,y,\xi)\Big|_{y=x}.
  \end{equation}
\end{theorem}

Here the terms with $\abs\alpha=j$ are grouped together; they have order
$m-j$. The symbol $\sigma_L(A)$ is the \emph{left} (or \emph{full})
\emph{symbol} of $A$. The proof occupies the rest of this section.

\subsection{Taylor expansion of the amplitude}

The key observation is the last identity of
\cref{euclidean:lem:kernel-identities}: multiplying the amplitude by
$(x-y)^\alpha$ appears to raise its weight, but actually lowers the order of
the operator by $\abs\alpha$.

\begin{lemma}\label{euclidean:lem:taylor}
  Let $a\in\Amp mw(\R^n)$ and $N\in\N$. Then
  \begin{equation*}
    \Op(a)=\sum_{j=0}^{N-1}\Op(a_j)+R_N,\qquad
    a_j(x,\xi)=\sum_{\abs\alpha=j}\frac{i^{\abs\alpha}}{\alpha!}
    D^\alpha_\xi D^\alpha_ya(x,y,\xi)\Big|_{y=x}\in S^{m-j}(\R^n;\R^n),
  \end{equation*}
  with $R_N\in\Psido{m-N}(\R^n)$; in fact $R_N=\Op(r_N)$ with
  $r_N\in\Amp{m-N}{w_+}(\R^n)$, $w_+=\max(w,0)$.
\end{lemma}
\begin{proof}
  Taylor's formula with integral remainder, applied to $a$ in the variable
  $y$ about $y=x$, gives
  \begin{equation*}
    a(x,y,\xi)=\sum_{\abs\alpha<N}\frac{(y-x)^\alpha}{\alpha!}(\pa^\alpha_ya)(x,x,\xi)
    +\sum_{\abs\alpha=N}(y-x)^\alpha\rho_\alpha(x,y,\xi),
  \end{equation*}
  \begin{equation*}
    \rho_\alpha(x,y,\xi)=\frac{N}{\alpha!}\int_0^1(1-t)^{N-1}
    (\pa^\alpha_ya)\bigl(x,x+t(y-x),\xi\bigr)\,dt.
  \end{equation*}
  We claim $\rho_\alpha\in\Amp m{w_+}$. Differentiating under the integral,
  $D^\mu_xD^\nu_yD^\gamma_\xi$ of the integrand is a linear combination, with
  coefficients that are polynomials in $t$ bounded on $[0,1]$, of
  derivatives of $\pa^\alpha_ya$ evaluated at $(x,x+t(y-x),\xi)$, and these
  are bounded by
  $C\ang{t(x-y)}^w\ang\xi^{m-\abs\gamma}\le C\ang{x-y}^{w_+}\ang\xi^{m-\abs\gamma}$,
  since $\ang{t(x-y)}\le\ang{x-y}$ and $\ang{t(x-y)}\ge 1$.
  Now apply \cref{euclidean:lem:kernel-identities}: for any
  $b\in\Amp{m}{w'}$,
  \begin{equation*}
    I\bigl((y-x)^\alpha b\bigr)=(-1)^{\abs\alpha}(x-y)^\alpha I(b)=I\bigl(D^\alpha_\xi b\bigr).
  \end{equation*}
  With $b=(\pa^\alpha_ya)(x,x,\xi)/\alpha!$, and $\pa_y^\alpha=i^{\abs\alpha}D^\alpha_y$,
  this shows that the term of order $\alpha$ in the Taylor sum contributes
  $\Op\bigl(\tfrac{i^{\abs\alpha}}{\alpha!}D^\alpha_\xi D^\alpha_ya(x,y,\xi)|_{y=x}\bigr)$.
  With $b=\rho_\alpha$ it shows that the remainder term contributes
  $\Op(D^\alpha_\xi\rho_\alpha)$, with
  $D^\alpha_\xi\rho_\alpha\in\Amp{m-N}{w_+}$.
\end{proof}

\subsection{Residual operators}

Applying Borel's lemma to the series in \cref{euclidean:lem:taylor} we will
find a symbol $b$ with $\Op(a)-\Op(b)\in\Psido{m-N}$ for all $N$. To finish
the proof of \cref{euclidean:thm:reduction} we need to identify the residual
operators. Note that \emph{a priori} an element of $\Psido{-\infty}(\R^n)$
is only known to be of the form $\Op(a_M)$ with $a_M\in\Amp{-M}{w_M}$ for each
$M$, where the weight $w_M$ may depend on $M$.

\begin{proposition}\label{euclidean:prop:residual}
  For a continuous linear operator $A:\cS(\R^n)\to\cSp(\R^n)$ with Schwartz
  kernel $K$ the following are equivalent:
  \begin{enumerate}[label=\textup{(\alph*)}]
  \item $A\in\Psido{-\infty}(\R^n)$;
  \item $K\in\CI(\R^{2n})$ and for all $\alpha,\beta\in\N_0^n$ and $N\in\N_0$
    \begin{equation}\label{euclidean:eq:residual-kernel}
      \abs{D^\alpha_xD^\beta_yK(x,y)}\le C_{\alpha\beta N}\ang{x-y}^{-N};
    \end{equation}
  \item $A=\Op(c)$ for some $c\in S^{-\infty}(\R^n;\R^n)$.
  \end{enumerate}
\end{proposition}
\begin{proof}
  (c)$\Rightarrow$(a) is immediate, since
  $S^{-\infty}(\R^n;\R^n)\subset\Amp M0(\R^n)$ for every $M$.

  (a)$\Rightarrow$(b). For each $M>n$ write $A=\Op(a_M)$ with
  $a_M\in\Amp{-M}{w_M}$. By \cref{euclidean:prop:osc-int}(i), $I(a_M)$ is a
  continuous, polynomially bounded function; two such functions which give
  the same integral against all $v\otimes u$ coincide, so $K=I(a_M)$ is a
  continuous function independent of $M$. By
  \cref{euclidean:lem:kernel-identities}, for all $\alpha,\beta,\gamma$,
  \begin{equation*}
    (x-y)^\gamma D^\alpha_xD^\beta_yK=I(b),\qquad
    b=(-D_\xi)^\gamma(\xi+D_x)^\alpha(D_y-\xi)^\beta a_M
    \in\Amp{-M+\abs\alpha+\abs\beta-\abs\gamma}{w_M},
  \end{equation*}
  where the powers of $\xi+D_x$ and $D_y-\xi$ are understood componentwise.
  Taking $\gamma=0$ and $M>n+\abs\alpha+\abs\beta$ shows that all
  distributional derivatives of $K$ are continuous functions, so
  $K\in\CI(\R^{2n})$; then $\abs{(x-y)^\gamma D^\alpha_xD^\beta_yK}\le C\ang{x-y}^{w_M}$
  for all $\gamma$. Since $\ang{x-y}^{2L}$ is a polynomial in $x-y$, this
  gives \eqref{euclidean:eq:residual-kernel}.

  (b)$\Rightarrow$(c). Set $g(x,z)=K(x,x-z)$. Since
  $D_xg=(D_xK+D_yK)(x,x-z)$ and $D_zg=-(D_yK)(x,x-z)$,
  \eqref{euclidean:eq:residual-kernel} gives
  $\abs{D^\alpha_xD^\beta_zg(x,z)}\le C_{\alpha\beta N}\ang z^{-N}$, so
  $\sup_{x,z}\abs{z^\gamma D^\alpha_xD^\beta_zg}<\infty$ for all
  $\alpha,\beta,\gamma$. Let
  \begin{equation*}
    c(x,\xi)=\int e^{-iz\cdot\xi}g(x,z)\,dz.
  \end{equation*}
  Then $\xi^\beta D^\gamma_\xi D^\alpha_xc$ is the Fourier transform in $z$ of
  $D^\beta_z\bigl((-z)^\gamma D^\alpha_xg\bigr)$, which is bounded in $L^1_z$
  uniformly in $x$; hence $\sup_{x,\xi}\abs{\xi^\beta D^\gamma_\xi D^\alpha_x c}<\infty$
  for all $\alpha,\beta,\gamma$, i.e.\ $c\in S^{-\infty}(\R^n;\R^n)$. By
  Fourier inversion $g(x,z)=\int e^{iz\cdot\xi}c(x,\xi)\,\dbar\xi$, so by
  \cref{euclidean:prop:osc-int}(i) the kernel of $\Op(c)$ is
  $\int e^{i(x-y)\cdot\xi}c(x,\xi)\,\dbar\xi=g(x,x-y)=K(x,y)$. Hence
  $A=\Op(c)$.
\end{proof}

\begin{proof}[Proof of \cref{euclidean:thm:reduction}]
  Let $A=\Op(a)$, $a\in\Amp mw$, and let $a_j\in S^{m-j}(\R^n;\R^n)$ be as in
  \cref{euclidean:lem:taylor}. By \cref{euclidean:prop:borel} choose
  $b\in S^m(\R^n;\R^n)$ with $b\sim\sum_ja_j$. Then for every $N$
  \begin{equation*}
    A-\Op(b)=R_N-\Op\Bigl(b-\sum_{j<N}a_j\Bigr)\in\Psido{m-N}(\R^n),
  \end{equation*}
  so $A-\Op(b)\in\Psido{-\infty}(\R^n)$ and, by
  \cref{euclidean:prop:residual}, $A-\Op(b)=\Op(c)$ with $c\in S^{-\infty}$.
  Thus $A=\Op(b+c)$ and $\sigma_L(A)=b+c$ satisfies
  \eqref{euclidean:eq:reduction}.

  For uniqueness, suppose $\Op(e)=0$ with $e\in S^m(\R^n;\R^n)$. By
  \eqref{euclidean:eq:left-quant},
  $\int e^{ix\cdot\xi}e(x,\xi)\hat u(\xi)\,\dbar\xi=0$ for all $u\in\cS$. For
  fixed $x$, $\hat u$ ranges over all of $\cS(\R^n)$, so the continuous,
  polynomially bounded function $\xi\mapsto e^{ix\cdot\xi}e(x,\xi)$ vanishes
  identically. Hence $e=0$.
\end{proof}

\begin{corollary}\label{euclidean:cor:reduction}
  \begin{enumerate}[label=\textup{(\roman*)}]
  \item $\Op:S^m(\R^n;\R^n)\to\Psido m(\R^n)$ is a bijection with inverse
    $\sigma_L$; in particular every $A\in\Psido m(\R^n)$ has an amplitude in
    $\Amp m0(\R^n)$, so the weights $w$ play no role in the definition of
    $\Psido m(\R^n)$.
  \item $\Psido{-\infty}(\R^n)=\Op\bigl(S^{-\infty}(\R^n;\R^n)\bigr)$, and
    for $A\in\Psido m(\R^n)$ and any $k\in\R$, $A\in\Psido{k}(\R^n)$ if and
    only if $\sigma_L(A)\in S^{k}(\R^n;\R^n)$.
  \item \textup{(Asymptotic summation of operators)} If
    $A_j\in\Psido{m-j}(\R^n)$, there exists $A\in\Psido m(\R^n)$, unique up to
    $\Psido{-\infty}(\R^n)$, with $A-\sum_{j<N}A_j\in\Psido{m-N}(\R^n)$ for
    all $N$; namely any $A$ with $\sigma_L(A)\sim\sum_j\sigma_L(A_j)$.
  \end{enumerate}
\end{corollary}
\begin{proof}
  (i) and (ii) are restatements of \cref{euclidean:thm:reduction} and
  \cref{euclidean:prop:residual}, using that $\sigma_L$ is linear. (iii)
  follows from (i), (ii) and \cref{euclidean:prop:borel}.
\end{proof}

% ======================================================================
\section{Composition and adjoints}\label{euclidean:sec:composition}

We can now prove the fundamental result of this chapter.

\begin{theorem}[Composition]\label{euclidean:thm:composition}
  The space $\Psido\infty(\R^n)$ is an order-filtered $*$-algebra of
  continuous operators on $\cS(\R^n)$ (and on $\cSp(\R^n)$):
  \begin{equation*}
    A\in\Psido m(\R^n)\Longrightarrow A^*\in\Psido m(\R^n),\qquad
    A\in\Psido m(\R^n),\ B\in\Psido{m'}(\R^n)\Longrightarrow AB\in\Psido{m+m'}(\R^n).
  \end{equation*}
  Moreover the left symbols satisfy
  \begin{gather}
    \sigma_L(A^*)(x,\xi)\sim\sum_\alpha\frac{i^{\abs\alpha}}{\alpha!}
    D^\alpha_\xi D^\alpha_x\overline{\sigma_L(A)(x,\xi)},\label{euclidean:eq:adjoint-symbol}\\
    \sigma_L(AB)(x,\xi)\sim\sum_\alpha\frac{i^{\abs\alpha}}{\alpha!}
    D^\alpha_\xi\sigma_L(A)(x,\xi)\,D^\alpha_x\sigma_L(B)(x,\xi).\label{euclidean:eq:composition-symbol}
  \end{gather}
\end{theorem}

Formally, \eqref{euclidean:eq:composition-symbol} reads
$\sigma_L(AB)=e^{iD_\eta\cdot D_y}\sigma_L(A)(x,\eta)\sigma_L(B)(y,\xi)\big|_{y=x,\eta=\xi}$.

\begin{proof}
  The statement about adjoints is \cref{euclidean:prop:adjoint}. By that
  proposition, the adjoint of $\Op(a)$, $a\in S^m(\R^n;\R^n)$, has amplitude
  $\overline{a(y,\xi)}$, so \cref{euclidean:thm:reduction} (with
  $D^\alpha_y\overline{a(y,\xi)}|_{y=x}=D^\alpha_x\overline{a(x,\xi)}$) gives
  \eqref{euclidean:eq:adjoint-symbol}.

  \emph{Right symbols.} Every $B\in\Psido{m'}(\R^n)$ can be written as
  $B=\Op_R(\sigma_R(B))$ with $\sigma_R(B)=\overline{\sigma_L(B^*)}\in S^{m'}$:
  indeed $B=(B^*)^*=\Op(\sigma_L(B^*))^*$, whose amplitude is
  $\overline{\sigma_L(B^*)(y,\xi)}$. Taking complex conjugates in
  \eqref{euclidean:eq:adjoint-symbol} applied to $B$, and using
  $\overline{D^\alpha f}=(-1)^{\abs\alpha}D^\alpha\bar f$ together with
  \cref{euclidean:lem:asymptotic-manip}(i),
  \begin{equation}\label{euclidean:eq:right-left}
    \sigma_R(B)\sim\sum_\beta\frac{(-i)^{\abs\beta}}{\beta!}
    D^\beta_\xi D^\beta_x\sigma_L(B).
  \end{equation}

  \emph{Composition.} Let $a=\sigma_L(A)$ and $b=\sigma_R(B)$. By
  \cref{euclidean:prop:left-right},
  $AB=\Op(a)\Op_R(b)=\Op\bigl(a(x,\xi)b(y,\xi)\bigr)$ with
  $a(x,\xi)b(y,\xi)\in\Amp{m+m'}0$, so $AB\in\Psido{m+m'}(\R^n)$, and
  \cref{euclidean:thm:reduction} gives
  \begin{equation*}
    \sigma_L(AB)\sim\sum_\alpha\frac{i^{\abs\alpha}}{\alpha!}
    D^\alpha_\xi\bigl[a\,D^\alpha_xb\bigr].
  \end{equation*}
  The map $f\mapsto D^\alpha_\xi[a\,D^\alpha_xf]$ lowers orders by
  $\abs\alpha$ and shifts them by $m$, so by
  \eqref{euclidean:eq:right-left} and
  \cref{euclidean:lem:asymptotic-manip}(i),(iii), with $c=\sigma_L(B)$,
  \begin{equation*}
    \sigma_L(AB)\sim\sum_\gamma\sum_{\alpha+\beta=\gamma}
    \frac{i^{\abs\alpha}(-i)^{\abs\beta}}{\alpha!\,\beta!}
    D^\alpha_\xi\bigl[a\,D^\beta_\xi D^\gamma_xc\bigr],
  \end{equation*}
  the terms with $\abs\gamma=l$ having order $m+m'-l$. By Leibniz' formula
  $D^\alpha_\xi[aD^\beta_\xi f]=\sum_{\mu\le\alpha}\binom\alpha\mu D^\mu_\xi a\,
  D^{\gamma-\mu}_\xi f$ when $\alpha+\beta=\gamma$, and using
  $\frac{\gamma!}{\alpha!\beta!}\binom\alpha\mu=\binom\gamma\mu\binom{\gamma-\mu}{\alpha-\mu}$,
  the coefficient of $\frac1{\gamma!}D^\mu_\xi a\,D^{\gamma-\mu}_\xi D^\gamma_xc$
  in the inner sum is
  \begin{equation*}
    \binom\gamma\mu i^{\abs\mu}\sum_{\delta\le\gamma-\mu}
    \binom{\gamma-\mu}\delta i^{\abs\delta}(-i)^{\abs{\gamma-\mu-\delta}}
    =\binom\gamma\mu i^{\abs\mu}\prod_{l=1}^n(i-i)^{\gamma_l-\mu_l},
  \end{equation*}
  which vanishes unless $\mu=\gamma$, when it is $i^{\abs\gamma}$. Hence the
  inner sum equals $\frac{i^{\abs\gamma}}{\gamma!}D^\gamma_\xi a\,D^\gamma_xc$,
  which is \eqref{euclidean:eq:composition-symbol}.

  The operators extend continuously to $\cSp(\R^n)$ by
  \eqref{euclidean:eq:S'-action}.
\end{proof}

The same computation, with the roles of the two expansions interchanged,
shows that left and right symbols are related by
\begin{equation}\label{euclidean:eq:left-right}
  \sigma_L\bigl(\Op_R(b)\bigr)\sim\sum_\alpha\frac{i^{\abs\alpha}}{\alpha!}
  D^\alpha_\xi D^\alpha_xb,
\end{equation}
which is just \cref{euclidean:thm:reduction} for the amplitude $b(y,\xi)$;
formally $\sigma_L=e^{iD_x\cdot D_\xi}\sigma_R$ and, by
\eqref{euclidean:eq:right-left}, $\sigma_R=e^{-iD_x\cdot D_\xi}\sigma_L$.
The map $\Op_R:S^m(\R^n;\R^n)\to\Psido m(\R^n)$ is a bijection with inverse
$\sigma_R$; see \cref{euclidean:ex:weyl} for a ``centrist'' alternative.

% ======================================================================
\section{The principal symbol}\label{euclidean:sec:principal}

By \eqref{euclidean:eq:left-right} and \eqref{euclidean:eq:right-left},
$\sigma_L(A)-\sigma_R(A)\in S^{m-1}(\R^n;\R^n)$ for $A\in\Psido m(\R^n)$: the
full symbol depends on the choice of quantization, but its leading part does
not. For this reason we consider the quotient spaces
\begin{equation*}
  S^{m}(\R^p;\R^n)/S^{m-1}(\R^p;\R^n),
\end{equation*}
and write $[a]$ for the class of $a\in S^m(\R^p;\R^n)$.

\begin{definition}\label{euclidean:def:principal-symbol}
  The \emph{principal symbol map}
  $\sigma_m:\Psido m(\R^n)\to S^m(\R^n;\R^n)/S^{m-1}(\R^n;\R^n)$ is
  $\sigma_m(A)=[\sigma_L(A)]=[\sigma_R(A)]$.
\end{definition}

Unlike $\sigma_L$ or $\sigma_R$, $\sigma_m$ depends on $m$: one needs to
know that the order is at most $m$ before it is defined. If $A=\Op(a)$ with
$a\in\Amp mw$, then $\sigma_m(A)=[a(x,x,\xi)]$ by
\cref{euclidean:thm:reduction}. For a differential operator
$P=\sum_{\abs\alpha\le m}p_\alpha D^\alpha\in\Diff m_{\mathrm b}(\R^n)$,
$\sigma_m(P)$ is represented by $\sum_{\abs\alpha=m}p_\alpha(x)\xi^\alpha$.

\begin{proposition}\label{euclidean:prop:symbol-sequence}
  For every $m\in\R$ the sequence
  \begin{equation*}
    0\longrightarrow\Psido{m-1}(\R^n)\hookrightarrow\Psido m(\R^n)
    \xrightarrow{\ \sigma_m\ }S^m(\R^n;\R^n)/S^{m-1}(\R^n;\R^n)\longrightarrow0
  \end{equation*}
  is exact. Moreover, for $A\in\Psido m(\R^n)$ and $B\in\Psido{m'}(\R^n)$,
  \begin{equation}\label{euclidean:eq:principal-mult}
    \sigma_{m+m'}(AB)=\sigma_m(A)\,\sigma_{m'}(B),\qquad
    \sigma_m(A^*)=\overline{\sigma_m(A)},
  \end{equation}
  and the commutator $[A,B]=AB-BA$ lies in $\Psido{m+m'-1}(\R^n)$ with
  \begin{equation}\label{euclidean:eq:commutator}
    \sigma_{m+m'-1}([A,B])=\frac1i\{\sigma_m(A),\sigma_{m'}(B)\},\qquad
    \{a,b\}=\sum_{j=1}^n\Bigl(\frac{\pa a}{\pa\xi_j}\frac{\pa b}{\pa x_j}
    -\frac{\pa a}{\pa x_j}\frac{\pa b}{\pa\xi_j}\Bigr).
  \end{equation}
\end{proposition}
\begin{proof}
  Surjectivity of $\sigma_m$ holds because $\sigma_m(\Op(a))=[a]$, and its
  null space is $\Psido{m-1}$ by \cref{euclidean:cor:reduction}(ii). All
  terms with $\abs\alpha\ge1$ in \eqref{euclidean:eq:composition-symbol} and
  \eqref{euclidean:eq:adjoint-symbol} have order at most $m+m'-1$, respectively
  $m-1$, which gives \eqref{euclidean:eq:principal-mult}. In particular
  $\sigma_{m+m'}([A,B])=0$, so $[A,B]\in\Psido{m+m'-1}$; the Poisson bracket
  is well defined on the classes (changing $a$ by an element of $S^{m-1}$
  changes $\{a,b\}$ by an element of $S^{m+m'-2}$). With $a=\sigma_L(A)$,
  $b=\sigma_L(B)$, the terms with $\abs\alpha=1$ in
  \eqref{euclidean:eq:composition-symbol} give
  \begin{equation*}
    \sigma_L([A,B])-i\sum_j\bigl(D_{\xi_j}a\,D_{x_j}b-D_{\xi_j}b\,D_{x_j}a\bigr)
    \in S^{m+m'-2},
  \end{equation*}
  and $iD_{\xi_j}a\,D_{x_j}b=-i\,\pa_{\xi_j}a\,\pa_{x_j}b$, which gives
  \eqref{euclidean:eq:commutator}.
\end{proof}

The simplicity of \eqref{euclidean:eq:principal-mult} compared with
\eqref{euclidean:eq:composition-symbol} is achieved at the expense of a great
loss of information. Still, many problems can be solved using it; it says
that the principal symbol gives a homomorphism from the filtered algebra
$\Psido0(\R^n)$ onto the \emph{commutative} algebra
$S^0(\R^n;\R^n)/S^{-1}(\R^n;\R^n)$.
